\documentclass[../main.tex]{subfiles}

\begin{document}

\chapter{度量化定理与仿紧致性}\label{chapter:ch06}

什么样的拓扑空间是可度量化的呢? 第 4 章中的 Urysohn 度量化定理朝着解决这个问题的方向迈出了第一步, 也是一大步, 给出了空间 $X$ 可以度量化的一个充分条件: $X$ 是正则的, 并且具有可数基. 但是只要能证得一个更强的结论, 数学家们就决不会满足于已有的这个定理. 现在, 人们有希望通过在 $X$ 上寻找一个使 $X$ 可度量化的充分必要条件, 即与可度量性等价的条件来强化这个定理.

我们知道, 在 Urysohn 度量化定理中正则性的假定是必要的, 而可数基的条件却不是. 于是一个很自然的想法就是用一些较弱的条件取代可数基条件. 寻找这样的条件需要一些技巧. 因为这个条件一方面要强到足以得出空间的可度量性, 另一方面又要弱到使每一个可度量化空间都能满足它. 但是, 难就难在如何发现正确的假设, 一旦发现了正确的假设, 离成功也就为期不远了.

这个条件终于被 J. Nagata 和 Y. Smirnov 各自独立找到了, 其中涉及到一个新概念, 即局部有限性概念. 空间 $X$ 的一个子集族 $\mathcal{A}$ 称为局部有限的, 指的是在 $X$ 的每一点处有一个邻域只和 $\mathcal{A}$ 中有限多个成员相交.

我们现在可以把基 $\mathcal{B}$ 可数这个条件说成 $\mathcal{B}$ 可以表为
\begin{align*}
\mathcal{B}=\bigcup_{n\in\mathbb{Z}_+}\mathcal{B}_n,
\end{align*}
其中每一个族 $\mathcal{B}_n$ 是有限的. 这是一个笨拙的表达方式, 但是它启发我们怎样去表达一个较弱的情形, Nagata-Smirnov 条件就是要求基 $\mathcal{B}$ 可以表示成以下形式
\begin{align*}
\mathcal{B}=\bigcup_{n\in\mathbb{Z}_+}\mathcal{B}_n,
\end{align*}
其中每一个族 $\mathcal{B}_n$ 是局部有限的. 这样的族 $\mathcal{B}$ 称为可数局部有限的. 出人意外的是这个条件加上正则性正好是 $X$ 可度量化的充分必要条件.

拓扑学中还有一个概念涉及局部有限性, 这就是所谓``仿紧致性''的概念, 它是紧致性概念的一个推广. 尽管这一概念的出现并不很久, 但已经在许多数学分支上被证明是有用的. 我们在这里引入它, 是为了给出 $X$ 可度量化的另一个充分必要条件, 即 $X$ 是可度量化的当且仅当 $X$ 既是仿紧致的又是局部可度量化的. 在第 42 节中将证明这个结论.

本章有些节之间是彼此独立的, 它们之间的依从关系如下表所示: 节 39 局部有限性引出节 40 Nagata-Smirnov 度量化定理与节 41 仿紧致性, 而后两者又共同引出节 42 Smirnov 度量化定理.

\subfile{Chapter-06/section-39}
\subfile{Chapter-06/section-40}
\subfile{Chapter-06/section-41}
\subfile{Chapter-06/section-42}

\end{document}